% Standalone kernel comparison accompanying
% "Physics-informed Bayesian optimization of high-order harmonic generation in gases"
\documentclass[11pt,a4paper]{article}
\usepackage[T1]{fontenc}
\usepackage[utf8]{inputenc}
\usepackage{lmodern}
\usepackage[margin=2.7cm]{geometry}
\usepackage{amsmath,amssymb,bm}
\usepackage{graphicx}
\usepackage{microtype}
\usepackage{caption}
\usepackage{hyperref}
\newcommand{\kernelfigurepath}{KERNELfigure2.png}
\begin{document}


\begin{center}
{\LARGE\bfseries Comparison of Gaussian Process Covariance Kernels for Bayesian Optimization of High-Order Harmonic Generation\par}

\vspace{1.1em}

{\normalsize
Dominika Maslarova$^{1,*}$, Mattys Pouyez$^{1}$,
Patrik Jansson$^{2}$, Jan Vabek$^{3,4}$,\\
Robin Weissenbilder$^{5}$, Linus Hagedorn$^{1,3}$,
Anne L'Huillier$^{5}$, and Tünde Fülöp$^{1}$\par}

\vspace{1em}

\begin{minipage}{0.95\textwidth}
\centering
\small

$^{1}$Department of Physics and Astronomy,
Chalmers University of Technology, Gothenburg, Sweden\par
\vspace{0.25em}

$^{2}$Department of Computer Science and Engineering,
Chalmers University of Technology and University of Gothenburg,
Gothenburg, Sweden\par
\vspace{0.25em}

$^{3}$Czech Technical University in Prague,
Faculty of Nuclear Sciences and Physical Engineering,
Prague, Czech Republic\par
\vspace{0.25em}

$^{4}$ELI Beamlines Facility,
The Extreme Light Infrastructure ERIC,
Doln\'{\i} B\v{r}e\v{z}any, Czech Republic\par
\vspace{0.25em}

$^{5}$Department of Physics,
Lund University, Lund, Sweden\par

\vspace{0.7em}

$^{*}$Corresponding author:
\texttt{dommas@chalmers.se}

\end{minipage}
\end{center}
\vspace{1.2em}
Here, we compare the performance of four covariance kernels for 2D Bayesian optimization using the LogEI acquisition function in the transformed optimization space $(\log(pL),\log(p/L))$, namely Mat\'ern $1/2$, Mat\'ern $3/2$, Mat\'ern $5/2$, and the radial basis function (RBF). These kernels encode differing assumptions regarding the smoothness of the Gaussian Process (GP) surrogate. The general Mat\'ern covariance function is given by
\begin{equation}
K_{\theta,\nu}(\mathbf{x},\mathbf{x}')
=s^2\frac{2^{1-\nu}}{\Gamma(\nu)}
\left(\sqrt{2\nu}\,R_{\boldsymbol{\ell}}\right)^\nu
K_\nu\!\left(\sqrt{2\nu}\,R_{\boldsymbol{\ell}}\right),
\label{eq:general_matern}
\end{equation}
where $R_{\boldsymbol{\ell}}$ is the scaled Euclidean distance between two points $\mathbf{x}=(x_1,\ldots,x_D)$ and $\mathbf{x'}=(x'_1,\ldots,x'_D)$ in the $D$-dimensional parameter space, given by
\begin{equation}
R_{\boldsymbol{\ell}}
=\sqrt{\sum_{j=1}^{D}\frac{(x_j-x'_j)^2}{\ell_j^2}},
\label{eq:matern_distance}
\end{equation}
where $\ell_j$ is the characteristic length scale associated with coordinate $x_j$, $s^2$ is the output variance, $\Gamma$ is the gamma function, and $K_\nu$ denotes the modified Bessel function of the second kind of order $\nu$. The parameter $\nu$ determines the smoothness of the GP prior. The Mat\'ern $1/2$, Mat\'ern $3/2$, and Mat\'ern $5/2$ kernels correspond to $\nu=1/2$, $3/2$, and $5/2$, respectively, whereas the RBF kernel is recovered in the limit $\nu\to\infty$, corresponding to an infinitely differentiable covariance function.

As shown in Figure~\ref{fig:kernel_fig1}a, the Mat\'ern $3/2$, Mat\'ern $5/2$, and RBF kernels exhibit nearly identical performance in terms of optimization convergence. Although the Mat\'ern $3/2$ kernel shows faster initial convergence, all three reach the same best-so-far observed conversion efficiency $\mathrm{CE}_{\mathrm{best}}$ by iteration 31. The Mat\'ern $1/2$ run was discontinued after iteration 10 because GP fitting failed, preventing a comparison of its long-term performance.

In terms of normalized root-mean-square error (NRMSE) and structural similarity index measure (SSIM) evaluated against the reference $100\times100$ grid scan (Figure~\ref{fig:kernel_fig1}b and c), the three completed runs achieve good surrogate reconstruction, although the RBF kernel exhibits somewhat lower predictive accuracy than the Mat\'ern $3/2$ and $5/2$ kernels. The latter two perform particularly similarly, with the Mat\'ern $5/2$ yielding slightly better final accuracy: NRMSE of approximately $0.039$ versus $0.045$, and SSIM of approximately $0.985$ versus $0.976$, respectively.
Consequently, the Mat\'ern $5/2$ kernel was retained for the main optimization studies as a robust choice, combining strong optimization convergence with the most accurate final surrogate reconstruction. While the Mat\'ern $3/2$ kernel converges somewhat faster initially, the differences between the two are small, and the RBF kernel remains a competitive alternative.

\begin{figure}[htbp]
    \centering
    \IfFileExists{\kernelfigurepath}{%
        \includegraphics[width=0.98\linewidth]{\kernelfigurepath}%
    }{%
        \fbox{\parbox[c][4cm][c]{0.88\linewidth}{\centering
        Place \texttt{KERNELfigure2.png} alongside this \texttt{.tex} file\\
        to display the original kernel-comparison figure.}}%
    }
    \caption{\textbf{Convergence and GP prediction accuracy for different kernels.}
    \textbf{a)} Evolution of the best-so-far CE as a function of the BO iteration.
    \textbf{b)} NRMSE between the GP prediction and the direct grid scan ($\mathrm{NRMSE}_{\mathrm{grid}}$) as a function of the BO iteration. The dotted gray horizontal line indicates the value of $0$, corresponding to perfect agreement with the direct grid scan.
    \textbf{c)} SSIM between the GP prediction grid and the direct grid scan ($\mathrm{SSIM}_{\mathrm{grid}}$) as a function of the BO iteration. The dotted gray horizontal line indicates the value of $1$, corresponding to perfect agreement with the direct grid scan. Different colors and marker styles distinguish the kernels.}
    \label{fig:kernel_fig1}
\end{figure}

\end{document}
